\documentclass[10pt,a4paper]{amsart}
\usepackage[margin=19mm]{geometry}
\usepackage{amsmath,amssymb,mathtools}
\usepackage[scr=rsfs,cal=euler]{mathalfa}
\usepackage{xfrac}
\usepackage[unicode,hidelinks,bookmarks=false]{hyperref}
\newcommand{\ie}{i.e.\ }
\newcommand{\cf}{cf.\ }
\newcommand{\resp}{respectively\ }
\newcommand{\Subsection}{\subsection}
\providecommand{\nobreakdash}{\nobreak\hbox{-}}
\setlength{\emergencystretch}{2em}
\multlinegap0pt
\usepackage{bm}
\usepackage{bbm}
\usepackage{dsfont}
\usepackage{xspace}
\usepackage{cancel}
\usepackage{mathtools}
\usepackage{amsthm}
\makeatletter
\@ifundefined{theorem}{\newtheorem{theorem}{Theorem}}{}
\@ifundefined{lemma}{\newtheorem{lemma}[theorem]{Lemma}}{}
\makeatother
\title[A sharp relative comparison inequality for conformal filling]{A sharp relative comparison inequality for conformal fillings of Poincar\'e--Einstein manifolds}
\author{Nan Wu}
\date{Source: arXiv:2607.13742v2 (CC BY 4.0)}
\begin{document}
\maketitle

\noindent\emph{Source and licence.} A sharp relative comparison inequality for conformal fillings of Poincar\'e--Einstein manifolds. Version arXiv:2607.13742v2 (CC BY 4.0), distributed under the Creative Commons Attribution 4.0 International licence, as recorded in the arXiv licence metadata for this exact version and shown on the abstract page https://arxiv.org/abs/2607.13742v2. Full rights notice: licenses/arxiv-2607.13742-CC-BY-4.0.txt.\par\smallskip

\noindent\textit{Editorial scope.} Counted unit: the Lemma whose source label is \texttt{lem:high-1D} in the article (source0.tex lines 856--879, source label \texttt{lem:high-1D}), delivered together with the article's own complete original proof of it (source0.tex lines 881--1004, one \texttt{proof} environment of 124 lines). No mathematics is written, repaired or re-derived: statement and proof are the article's own text line for line. Required context retained uncounted: eq:beta-via-H (lines 549--551); lem:beta-positive (lines 611--620); eq: def of LN PN 1D (lines 656--658); eq:H-N (lines 841--843); eq:H-ODE (lines 845--851); eq:high-f (lines 853--855). Every definition, display and label of those lines is retained unchanged. Omitted from this package and not redistributed: 1--548, 552--610, 621--655, 659--840, 844--844, 852--852, 880--880, 1005--1302. Nothing inside a retained excerpt is omitted or altered: every retained body is delivered exactly as the article's lines are written, so no omission receipt of any kind accompanies this package. Citations: the retained excerpts carry no citation command at all, so no bibliography entry of the article is transcribed and no reference is redistributed. Reference adaptation: none was needed and none was made; every reference inside the delivered unit resolves inside it, so no unresolved reference or citation is produced. Printed slips kept literal: no printed word, symbol, hypothesis or number of the counted statement or of its proof was altered. Layout and class: the article's own class is \texttt{amsart} with packages including fontenc, inputenc, amsmath, amssymb, amsthm, mathtools, mathrsfs, enumitem, hyperref; the wrapper is the lane's pinned \texttt{amsart} editorial wrapper at 10pt on A4, which restates the theorem-like environments the retained text needs and adds the article's own macro definitions that the retained text needs (none), with extra wrapper packages (bm, bbm, dsfont, xspace, cancel, mathtools). The delivered numbering is the unit's own. Assets: no retained span contains a figure environment, a graphics-inclusion command, a PDF-page inclusion, a table or a TikZ picture; no image file is copied into this package, the package declares no asset dependency and no image file is redistributed with it. Licence: version arXiv:2607.13742v2 of this article is distributed under the Creative Commons Attribution 4.0 International licence, as recorded in the arXiv licence metadata for this exact version and shown on the abstract page https://arxiv.org/abs/2607.13742v2. A full rights notice is delivered at licenses/arxiv-2607.13742-CC-BY-4.0.txt. Characters: every retained excerpt is pure ASCII, so the delivered engine, XeLaTeX with its default Latin Modern text font, reports no missing character.
\begin{equation}\label{eq:beta-via-H}
 \beta_N(r)=C_N\frac{(1+r^2)H_N(r)-2r/N}{1-r^2}.
\end{equation}

\begin{lemma}\label{lem:beta-positive}
For every $N\ge3$, the function $\beta_N(r)$ satisfies
\[
 \beta_N(r)>0\qquad \text{ for}\quad  0\le r\le1.
\]
Consequently,
\begin{equation}\label{eq:reverse-bridge-general}
 A\ge C_NV.
\end{equation}
\end{lemma}

\begin{equation}\label{eq: def of LN PN 1D}
L_N f=-\left(1-r^2\right) f^{\prime \prime}+N r f^{\prime}, \quad P_N[f]=f^{\prime}+\frac{2 r}{N} f^{\prime \prime},
\end{equation}

\begin{equation}\label{eq:H-N}
 H_N(r)=\frac{K_N(r)}{(1-r^2)^{N/2}},
\end{equation}

\begin{equation}\label{eq:H-ODE}
 H_N'(r)=\frac{NrH_N(r)-1}{1-r^2},
 \qquad
 H_N(0)=\frac1{C_N},
 \qquad
 \lim_{r\to1}H_N(r)=\frac1N.
\end{equation}

\begin{equation}\label{eq:high-f}
 f_N(r)=\frac{C_N}{N}\left(\frac1{C_N}-H_N(r)\right).
\end{equation}

\begin{lemma}\label{lem:high-1D}
For $N\ge5$, let $H_N$ and $f_N$ be the functions as in \eqref{eq:H-N} and \eqref{eq:high-f}. Set
\begin{equation}\label{eq:calP-def}
 P_N(r)=P_N\left[f_N\right](r)=f_N'(r)+\frac{2r}{N}f_N''(r).
\end{equation}
Then,
\begin{equation}\label{eq:high-f-bounds}
 0\le f_N(r)\le f_N(1)=\frac1N\left(1-\frac{C_N}{N}\right).
\end{equation}
Moreover, 
\begin{equation}\label{eq:high-Lf}
 L_Nf_N=\beta_N,
\end{equation}
where 
\[
 L_Nf=-(1-r^2)f''+Nrf',
\]
as in \eqref{eq: def of LN PN 1D}.
Furthermore, 
\begin{equation}\label{eq:high-pointwise}
\beta_N(r)^2\le2C_NP_N(r)\qquad \text{ for } 0\le r\le1.
\end{equation}
In particular, $P_N(r)>0$ on $[0,1]$.
\end{lemma}

\begin{proof}
We first prove \eqref{eq:high-f-bounds}. The lower bound estimate for $H_N$ follows from 
\[K_N(r)=\frac{1}{2} \int_0^{1-r^2} u^{N / 2-1}(1-u)^{-1 / 2} d u \geq \frac{1}{2} \int_0^{1-r^2} u^{N / 2-1} d u=\frac{\left(1-r^2\right)^{N / 2}}{N}.\]
Then, $H_N \ge 1/N$. For $r>0$,
\[
K_N(r) \leq \frac{1}{r} \int_r^1 t\left(1-t^2\right)^{(N-2) / 2} d t=\frac{\left(1-r^2\right)^{N / 2}}{N r}.
\]
By \eqref{eq:H-ODE},
\[
 H_N'(r)=\frac{NrH_N(r)-1}{1-r^2}\le0
\]
on $(0,1)$, and then $H_N(r)\le H_N(0)=1/C_N$.  Hence, $ 0\le f_N(r)\le f_N(1),$ which implies \eqref{eq:high-f-bounds}.

Differentiating \eqref{eq:high-f} and using \eqref{eq:H-ODE} give
\begin{equation}\label{eq:fprime-high-proof}
 f_N'(r)=\frac{C_N}{N}\frac{1-NrH_N(r)}{1-r^2}.
\end{equation}
Substituting \eqref{eq:H-ODE} into the derivative of \eqref{eq:fprime-high-proof} yields
\[
 L_Nf_N=-(1-r^2)f_N''+Nrf_N'
 =C_N\frac{(1+r^2)H_N(r)-2r/N}{1-r^2}.
\]
By \eqref{eq:beta-via-H}, the last expression is exactly $\beta_N$. This proves \eqref{eq:high-Lf}.

It remains to prove the pointwise inequality \eqref{eq:high-pointwise}.  A direct calculation from \eqref{eq:H-ODE} and \eqref{eq:fprime-high-proof} gives 
\begin{equation}\label{eq:calP-explicit}
 \frac{P_N(r)}{C_N}
 =\frac{N+(N+4)r^2-N(N+2)r(1+r^2)H_N(r)}{N^2(1-r^2)^2}.
\end{equation}
Using 
\[
\frac{\beta_N(r)}{C_N}=\frac{\left(1+r^2\right) H_N(r)-2 r / N}{1-r^2},
\]
which is \eqref{eq:beta-via-H}, together with \eqref{eq:calP-explicit}, we know that
$\beta_N^2\le2C_NP_N$ is equivalent to
\begin{equation}\label{eq:H-upper-needed}
 \left((1+r^2)H_N(r)+r\right)^2
 \le
 r^2+\frac{2(1+r^2)}{N}+\frac{4r^2}{N^2}.
\end{equation}
Define
\begin{equation}\label{eq:Hplus}
 H_N^+(r)=
 \frac{-r+\left(r^2+\frac{2(1+r^2)}{N}+\frac{4r^2}{N^2}\right)^{1/2}}
 {1+r^2}.
\end{equation}
Then, \eqref{eq:H-upper-needed} is equivalent to
\begin{equation}\label{eq:HleHplus}
 H_N(r)\le H_N^+(r).
\end{equation}

We now prove this comparison inequality.

Set
\[
 B_N(r)=2N+(N^2+2N+4)r^2.
\]
Then
\[
 H_N^+(r)=\frac{-r+N^{-1}\sqrt{B_N(r)}}{1+r^2}.
\]
Direct differentiation gives
\begin{equation}\label{eq:Hplus-residual}
 (H_N^+)'-\frac{NrH_N^+-1}{1-r^2}
 =-\frac{r\,Z_N(r)}{N(1-r^2)(1+r^2)^2\sqrt{B_N(r)}},
\end{equation}
where
\begin{align}
 \mathcal A_N(r)
 &=(N^3+N^2+2N-4)r^4\label{eq:AcalN}\\ 
 &\qquad+(N^3+6N^2+4N+8)r^2
 +(N^2+2N-4),\\
 Z_N(r)
 &=\mathcal A_N(r)
 -r\sqrt{B_N(r)}\,\bigl(N^2(1+r^2)+4N\bigr).\label{eq:XiN}
\end{align}
We claim that $Z_N(r) \geq 0$ for $0 \leq r \leq 1$.
The polynomial $\mathcal A_N$ is strictly positive on $[0,1]$ for $N\ge5$.  Moreover $N^2(1+r^2)+4N>0$. Squaring the two nonnegative terms in \eqref{eq:XiN}, we obtain 
\begin{equation}\label{eq:Sidentity}
 A_N(r)^2
 -r^2B_N(r)\bigl(N^2(1+r^2)+4N\bigr)^2
=(1-r^2)^2\mathcal S_N(r),
\end{equation}
where
\begin{align}
 \mathcal S_N(r)
 &=(N^4-4N^3-4N^2-16N+16)r^4 \notag\\
 &\quad +(2N^4-24N^2-32N-32)r^2
 \\
 &\qquad+N^4+4N^3-4N^2-16N+16.\label{eq:S-explicit}
\end{align}
For $N\ge6$, all three coefficients in \eqref{eq:S-explicit} are positive.  For $N=5$,
\[
 \mathcal S_5(r)=-39r^4+458r^2+961,
\]
which is also positive for $0\le r\le1$.  Hence $\mathcal S_N(r)\ge0$ for all $N\ge5$.  Since $\mathcal A_N\ge0$ and the second term in \eqref{eq:XiN} is also nonnegative, \eqref{eq:Sidentity} implies
\[
 \mathcal A_N(r)
 \ge r\sqrt{B_N(r)}\,\bigl(N^2(1+r^2)+4N\bigr),
\]
and therefore $Z_N(r)\ge0$.  Thus by \eqref{eq:Hplus-residual},
\begin{equation}\label{eq:Hplus-subsolution}
 (H_N^+)'-\frac{NrH_N^+-1}{1-r^2}\le0.
\end{equation}
Finally set $z=H_N^+-H_N$.  Combining \eqref{eq:H-ODE} and \eqref{eq:Hplus-subsolution} gives
\[
 z'-\frac{Nr}{1-r^2}z\le0.
\]
Multiplying by the integrating factor $(1-r^2)^{N/2}$, we obtain
\[
 \frac{d}{dr}\Big((1-r^2)^{N/2}z(r)\Big)\le0.
\]
Both $H_N$ and $H_N^+$ extend continuously to $r=1$ and satisfy
\[
 H_N(1)=H_N^+(1)=\frac1N.
\]
Thus, for $0<r<R<1$, integrate from $r$ to $R$ and then let $R\to 1$.  Since $z$ is bounded at $1$,
\[
 \lim_{R\to 1}(1-R^2)^{N/2}z(R)=0.
\]
Thus $(1-r^2)^{N/2}z(r)\ge0$, hence $z(r)\ge0$ on $[0,1)$. $z(1)\ge 0$ follows by continuity.  This proves \eqref{eq:HleHplus}, hence \eqref{eq:H-upper-needed}, and therefore \eqref{eq:high-pointwise}.  

Finally, Lemma \ref{lem:beta-positive} gives $\beta_N>0$ on $[0,1]$.  Therefore \eqref{eq:high-pointwise} implies $P_N>0$ on $[0,1]$.
\end{proof}

\end{document}
